% Physical page 2 of the concise study, which carries this sheet and nothing
% else. The dates, relation labels, alternatives and unresolved states are the
% generated projection of research/chronology.toml, and stand exactly as the
% reviewed appendix of the expansive study prints them; the explanatory rows
% condense that appendix's. The generated definitions are imported here, once,
% because the concise entrypoint's own preamble does not carry them.
% Generated. Do not edit.
% Deterministic projection of research/chronology.toml.
% schema: 3
% document: liturgy/roman-rite/1962/propers/temporal/59-nineteenth-after-pentecost
% generated-by: tools/tpt proper-chronology annotations
\newcommand{\chronologyannotationclaim}[6]{#6}
\newcommand{\chronologyannotationcomparisonclaim}[8]{#8}
\newcommand{\chronologyannotationreach}[3]{}
\newcommand{\chronologyannotationgroup}[3]{#3}
\newcommand{\chronologyannotationcomparisongroup}[4]{#4}
\newcommand{\chronologyannotation}[1]{%
  \ifcsname triptychchronologyannotation@#1\endcsname
    \csname triptychchronologyannotation@#1\endcsname
  \else
    \PackageError{triptych}{No chronology annotation for #1}{}%
  \fi}
% comparison-dependency: src/claude/liturgy/roman-rite/1962/propers/temporal/59-nineteenth-after-pentecost/research/chronology-profile-comparisons.toml sha256=1c48dd33c95c1d45ba53f672e8a3b9cfe81060ac70930f37ac3d8e912e365f08
% introit: Introit
\expandafter\def\csname triptychchronologyannotation@introit\endcsname{%
  \chronologyannotationgroup{composition}{preferred}{\textbf{Composition}: \chronologyannotationreach{Ps.77.1}{Ps}{true}\chronologyannotationclaim{critical.psalms.latest-composition-boundary}{composition}{catholic-critical-v1}{preferred}{before the Maccabean period, around 165 B.C.; no individual psalm can be dated securely}{Before c. 165 B.C.}}%
}
% epistle: Epistle
\expandafter\def\csname triptychchronologyannotation@epistle\endcsname{%
  \chronologyannotationgroup{composition}{disputed}{\textbf{Composition} -- disputed: \chronologyannotationreach{Eph.4.23}{Eph}{true}\chronologyannotationreach{Eph.4.24}{Eph}{true}\chronologyannotationreach{Eph.4.25}{Eph}{true}\chronologyannotationreach{Eph.4.26}{Eph}{true}\chronologyannotationreach{Eph.4.27}{Eph}{true}\chronologyannotationreach{Eph.4.28}{Eph}{true}\chronologyannotationclaim{composition.epistle-to-the-ephesians}{composition}{catholic-traditional-v1}{disputed}{a period between 58 and 63}{A.D. 58--63}; \chronologyannotationreach{Eph.4.23}{Eph}{true}\chronologyannotationreach{Eph.4.24}{Eph}{true}\chronologyannotationreach{Eph.4.25}{Eph}{true}\chronologyannotationreach{Eph.4.26}{Eph}{true}\chronologyannotationreach{Eph.4.27}{Eph}{true}\chronologyannotationreach{Eph.4.28}{Eph}{true}\chronologyannotationclaim{composition.epistle-to-the-ephesians}{composition}{catholic-traditional-v1}{disputed}{(Philemon; Colossians; Ephesians; Philippians), 61}{A.D. 61}.}%
  \space \chronologyannotationcomparisongroup{catholic-critical-v1}{composition}{composition-only}{\textbf{Composition (Catholic critical chronology, for comparison)} -- disputed: \chronologyannotationreach{Eph.4.23}{Eph}{true}\chronologyannotationreach{Eph.4.24}{Eph}{true}\chronologyannotationreach{Eph.4.25}{Eph}{true}\chronologyannotationreach{Eph.4.26}{Eph}{true}\chronologyannotationreach{Eph.4.27}{Eph}{true}\chronologyannotationreach{Eph.4.28}{Eph}{true}\chronologyannotationcomparisonclaim{critical.ephesians-later-disciple}{composition}{catholic-critical-v1}{catholic-critical-v1}{disputed}{passage.united-states-conference-of-catholic-bishops.new-american-bible-revised-edition.english-usccb-web-2026-09-28.ephesians-introduction}{around A.D. 80--100}{Ephesians under the later-disciple hypothesis (approximate), c. A.D. 80--100}.}%
}
% gradual: Gradual
\expandafter\def\csname triptychchronologyannotation@gradual\endcsname{%
  \chronologyannotationgroup{traditional-attribution}{disputed}{\textbf{Traditional attribution} -- disputed: \chronologyannotationreach{Ps.140.2}{Ps.137 Ps.140}{true}\chronologyannotationclaim{israel.monarchy.david-traditional-era}{traditional-attribution}{catholic-traditional-v1}{disputed}{reigned from 1055 to 1015 B.C.}{David (reign in the usual chronology), B.C. 1055--1015}.}%
  \space \chronologyannotationgroup{composition}{preferred}{\textbf{Composition}: \chronologyannotationreach{Ps.140.2}{Ps}{true}\chronologyannotationclaim{critical.psalms.latest-composition-boundary}{composition}{catholic-critical-v1}{preferred}{before the Maccabean period, around 165 B.C.; no individual psalm can be dated securely}{Before c. 165 B.C.}}%
}
% alleluia: Alleluia
\expandafter\def\csname triptychchronologyannotation@alleluia\endcsname{%
  \chronologyannotationgroup{composition}{preferred}{\textbf{Composition}: \chronologyannotationreach{Ps.104.1}{Ps}{true}\chronologyannotationclaim{critical.psalms.latest-composition-boundary}{composition}{catholic-critical-v1}{preferred}{before the Maccabean period, around 165 B.C.; no individual psalm can be dated securely}{Before c. 165 B.C.}}%
}
% gospel: Gospel
\expandafter\def\csname triptychchronologyannotation@gospel\endcsname{%
  \chronologyannotationgroup{narrated-event}{preferred}{\textbf{Event}: \chronologyannotationreach{Matt.22.1}{Matt.22.1-14}{false}\chronologyannotationreach{Matt.22.10}{Matt.22.1-14}{false}\chronologyannotationreach{Matt.22.11}{Matt.22.1-14}{false}\chronologyannotationreach{Matt.22.12}{Matt.22.1-14}{false}\chronologyannotationreach{Matt.22.13}{Matt.22.1-14}{false}\chronologyannotationreach{Matt.22.14}{Matt.22.1-14}{false}\chronologyannotationreach{Matt.22.2}{Matt.22.1-14}{false}\chronologyannotationreach{Matt.22.3}{Matt.22.1-14}{false}\chronologyannotationreach{Matt.22.4}{Matt.22.1-14}{false}\chronologyannotationreach{Matt.22.5}{Matt.22.1-14}{false}\chronologyannotationreach{Matt.22.6}{Matt.22.1-14}{false}\chronologyannotationreach{Matt.22.7}{Matt.22.1-14}{false}\chronologyannotationreach{Matt.22.8}{Matt.22.1-14}{false}\chronologyannotationreach{Matt.22.9}{Matt.22.1-14}{false}\chronologyannotationclaim{life-of-christ.parable-of-the-marriage-feast}{narrated-event}{catholic-traditional-v1}{preferred}{derived A.D. 29 from February, A.U.C. 782- Passover, 782}{A.D. 29 (derived)}.}%
  \space \chronologyannotationgroup{composition}{disputed}{\textbf{Composition} -- disputed: \chronologyannotationreach{Matt.22.1}{Matt}{true}\chronologyannotationreach{Matt.22.10}{Matt}{true}\chronologyannotationreach{Matt.22.11}{Matt}{true}\chronologyannotationreach{Matt.22.12}{Matt}{true}\chronologyannotationreach{Matt.22.13}{Matt}{true}\chronologyannotationreach{Matt.22.14}{Matt}{true}\chronologyannotationreach{Matt.22.2}{Matt}{true}\chronologyannotationreach{Matt.22.3}{Matt}{true}\chronologyannotationreach{Matt.22.4}{Matt}{true}\chronologyannotationreach{Matt.22.5}{Matt}{true}\chronologyannotationreach{Matt.22.6}{Matt}{true}\chronologyannotationreach{Matt.22.7}{Matt}{true}\chronologyannotationreach{Matt.22.8}{Matt}{true}\chronologyannotationreach{Matt.22.9}{Matt}{true}\chronologyannotationclaim{composition.gospel-of-matthew}{composition}{catholic-traditional-v1}{disputed}{about A.D. 38-45}{c. A.D. 38--45}; \chronologyannotationreach{Matt.22.1}{Matt}{true}\chronologyannotationreach{Matt.22.10}{Matt}{true}\chronologyannotationreach{Matt.22.11}{Matt}{true}\chronologyannotationreach{Matt.22.12}{Matt}{true}\chronologyannotationreach{Matt.22.13}{Matt}{true}\chronologyannotationreach{Matt.22.14}{Matt}{true}\chronologyannotationreach{Matt.22.2}{Matt}{true}\chronologyannotationreach{Matt.22.3}{Matt}{true}\chronologyannotationreach{Matt.22.4}{Matt}{true}\chronologyannotationreach{Matt.22.5}{Matt}{true}\chronologyannotationreach{Matt.22.6}{Matt}{true}\chronologyannotationreach{Matt.22.7}{Matt}{true}\chronologyannotationreach{Matt.22.8}{Matt}{true}\chronologyannotationreach{Matt.22.9}{Matt}{true}\chronologyannotationclaim{composition.gospel-of-matthew}{composition}{catholic-traditional-v1}{disputed}{about the year 40-42}{c. A.D. 40--42}; \chronologyannotationreach{Matt.22.1}{Matt}{true}\chronologyannotationreach{Matt.22.10}{Matt}{true}\chronologyannotationreach{Matt.22.11}{Matt}{true}\chronologyannotationreach{Matt.22.12}{Matt}{true}\chronologyannotationreach{Matt.22.13}{Matt}{true}\chronologyannotationreach{Matt.22.14}{Matt}{true}\chronologyannotationreach{Matt.22.2}{Matt}{true}\chronologyannotationreach{Matt.22.3}{Matt}{true}\chronologyannotationreach{Matt.22.4}{Matt}{true}\chronologyannotationreach{Matt.22.5}{Matt}{true}\chronologyannotationreach{Matt.22.6}{Matt}{true}\chronologyannotationreach{Matt.22.7}{Matt}{true}\chronologyannotationreach{Matt.22.8}{Matt}{true}\chronologyannotationreach{Matt.22.9}{Matt}{true}\chronologyannotationclaim{composition.gospel-of-matthew}{composition}{catholic-traditional-v1}{disputed}{the years 40-45}{A.D. 40--45}; \chronologyannotationreach{Matt.22.1}{Matt}{true}\chronologyannotationreach{Matt.22.10}{Matt}{true}\chronologyannotationreach{Matt.22.11}{Matt}{true}\chronologyannotationreach{Matt.22.12}{Matt}{true}\chronologyannotationreach{Matt.22.13}{Matt}{true}\chronologyannotationreach{Matt.22.14}{Matt}{true}\chronologyannotationreach{Matt.22.2}{Matt}{true}\chronologyannotationreach{Matt.22.3}{Matt}{true}\chronologyannotationreach{Matt.22.4}{Matt}{true}\chronologyannotationreach{Matt.22.5}{Matt}{true}\chronologyannotationreach{Matt.22.6}{Matt}{true}\chronologyannotationreach{Matt.22.7}{Matt}{true}\chronologyannotationreach{Matt.22.8}{Matt}{true}\chronologyannotationreach{Matt.22.9}{Matt}{true}\chronologyannotationclaim{composition.gospel-of-matthew}{composition}{catholic-traditional-v1}{disputed}{about the year 60-68}{c. A.D. 60--68}; \chronologyannotationreach{Matt.22.1}{Matt}{true}\chronologyannotationreach{Matt.22.10}{Matt}{true}\chronologyannotationreach{Matt.22.11}{Matt}{true}\chronologyannotationreach{Matt.22.12}{Matt}{true}\chronologyannotationreach{Matt.22.13}{Matt}{true}\chronologyannotationreach{Matt.22.14}{Matt}{true}\chronologyannotationreach{Matt.22.2}{Matt}{true}\chronologyannotationreach{Matt.22.3}{Matt}{true}\chronologyannotationreach{Matt.22.4}{Matt}{true}\chronologyannotationreach{Matt.22.5}{Matt}{true}\chronologyannotationreach{Matt.22.6}{Matt}{true}\chronologyannotationreach{Matt.22.7}{Matt}{true}\chronologyannotationreach{Matt.22.8}{Matt}{true}\chronologyannotationreach{Matt.22.9}{Matt}{true}\chronologyannotationclaim{composition.gospel-of-matthew}{composition}{catholic-traditional-v1}{disputed}{about the years 64-67}{c. A.D. 64--67}; \chronologyannotationreach{Matt.22.1}{Matt}{true}\chronologyannotationreach{Matt.22.10}{Matt}{true}\chronologyannotationreach{Matt.22.11}{Matt}{true}\chronologyannotationreach{Matt.22.12}{Matt}{true}\chronologyannotationreach{Matt.22.13}{Matt}{true}\chronologyannotationreach{Matt.22.14}{Matt}{true}\chronologyannotationreach{Matt.22.2}{Matt}{true}\chronologyannotationreach{Matt.22.3}{Matt}{true}\chronologyannotationreach{Matt.22.4}{Matt}{true}\chronologyannotationreach{Matt.22.5}{Matt}{true}\chronologyannotationreach{Matt.22.6}{Matt}{true}\chronologyannotationreach{Matt.22.7}{Matt}{true}\chronologyannotationreach{Matt.22.8}{Matt}{true}\chronologyannotationreach{Matt.22.9}{Matt}{true}\chronologyannotationclaim{composition.gospel-of-matthew}{composition}{catholic-traditional-v1}{disputed}{about the year 50}{c. A.D. 50}.}%
  \space \chronologyannotationcomparisongroup{catholic-critical-v1}{composition}{composition-only}{\textbf{Composition (Catholic critical chronology, for comparison)}: \chronologyannotationreach{Matt.22.1}{Matt}{true}\chronologyannotationreach{Matt.22.10}{Matt}{true}\chronologyannotationreach{Matt.22.11}{Matt}{true}\chronologyannotationreach{Matt.22.12}{Matt}{true}\chronologyannotationreach{Matt.22.13}{Matt}{true}\chronologyannotationreach{Matt.22.14}{Matt}{true}\chronologyannotationreach{Matt.22.2}{Matt}{true}\chronologyannotationreach{Matt.22.3}{Matt}{true}\chronologyannotationreach{Matt.22.4}{Matt}{true}\chronologyannotationreach{Matt.22.5}{Matt}{true}\chronologyannotationreach{Matt.22.6}{Matt}{true}\chronologyannotationreach{Matt.22.7}{Matt}{true}\chronologyannotationreach{Matt.22.8}{Matt}{true}\chronologyannotationreach{Matt.22.9}{Matt}{true}\chronologyannotationcomparisonclaim{critical.gospel-of-matthew}{composition}{catholic-critical-v1}{catholic-critical-v1}{preferred}{passage.united-states-conference-of-catholic-bishops.new-american-bible-revised-edition.english-usccb-web-2026-09-21.matthew-introduction}{post-A.D. 70 date}{The Greek Gospel of Matthew in the NABRE introduction, post-A.D. 70 date}.}%
}
% offertory: Offertory
\expandafter\def\csname triptychchronologyannotation@offertory\endcsname{%
  \chronologyannotationgroup{traditional-attribution}{disputed}{\textbf{Traditional attribution} -- disputed: \chronologyannotationreach{Ps.137.7}{Ps.137 Ps.140}{true}\chronologyannotationclaim{israel.monarchy.david-traditional-era}{traditional-attribution}{catholic-traditional-v1}{disputed}{reigned from 1055 to 1015 B.C.}{David (reign in the usual chronology), B.C. 1055--1015}.}%
  \space \chronologyannotationgroup{composition}{preferred}{\textbf{Composition}: \chronologyannotationreach{Ps.137.7}{Ps}{true}\chronologyannotationclaim{critical.psalms.latest-composition-boundary}{composition}{catholic-critical-v1}{preferred}{before the Maccabean period, around 165 B.C.; no individual psalm can be dated securely}{Before c. 165 B.C.}}%
}
% communion: Communion
\expandafter\def\csname triptychchronologyannotation@communion\endcsname{%
  \chronologyannotationgroup{composition}{preferred}{\textbf{Composition}: \chronologyannotationreach{Ps.118.4}{Ps}{true}\chronologyannotationreach{Ps.118.5}{Ps}{true}\chronologyannotationclaim{critical.psalms.latest-composition-boundary}{composition}{catholic-critical-v1}{preferred}{before the Maccabean period, around 165 B.C.; no individual psalm can be dated securely}{Before c. 165 B.C.}}%
}


\section{Scriptural Date and Location}

\zlabel{triptych:concise:chronology:start}
\begin{concisedossiertable}
Introit (verse) & Ps 77:1 (Heb.\ 78) & Asaph, by the title, to Israel as ``my
people''; a history from Egypt to Sion and David (vv.~68--70) &
\chronodate{introit}{\chronologyannotation{introit}}\\
\cmidrule(lr){1-4}
\dossierprose{Titled in the Vulgate \latin{Intellectus Asaph}, ``Understanding
for Asaph'', which the Missal's verse leaves out; the title gives no occasion or
date. The displayed date is the bound that the \work{New American Bible,
Revised Edition} sets for the whole Psalter: its introduction to the Psalms
holds that no psalm is as late as the Maccabean period and that none can be
dated securely. The antiphon stands in no verse of the Vulgate and has no
dossier of its own.}
\midrule
Alleluia & Ps 104:1 (Heb.\ 105) & Israel, bidden to declare God's deeds among
the Gentiles; sung before the ark at Jerusalem (1 Par 15:3; 16:7--8) &
\chronodate{alleluia}{\chronologyannotation{alleluia}}\\
\cmidrule(lr){1-4}
\dossierprose{Titled \latin{Alleluia} in the Vulgate, which the Missal's verse
leaves out. Paralipomenon sets these words in the praise Asaph sang on the day
the ark was set in David's tent (1 Par 16:1, 7); the displayed date is the
Psalter's bound, not that day's.}
\midrule
Communion & Ps 118:4--5 (Heb.\ 119) & The psalmist, ``a sojourner on the
earth'' (v.~19), singing the law ``in the place of my pilgrimage'' (v.~54) &
\chronodate{communion}{\chronologyannotation{communion}}\\
\cmidrule(lr){1-4}
\dossierprose{Titled \latin{Alleluia}, with no author named; the Communion's
verses stand in its first alphabetic strophe, \latin{Aleph}. The displayed date
is the Psalter's bound.}
\midrule
Offertory & Ps 137:7 (Heb.\ 138) & David, by the title; praise ``in the sight
of the angels'' and worship ``towards thy holy temple'' (vv.~1--2) &
\chronodate{offertory}{\chronologyannotation{offertory}}\\
\cmidrule(lr){1-4}
\dossierprose{Titled \latin{Ipsi David}, ``For David himself'', without an
occasion. The reign shown is David's in the ``usual chronology'' of the
\work{Catholic Encyclopedia}'s ``David, King'' (vol.~4, 1908), which reports a later dating
beside it: the title's reference point, not the psalm's date. The composition
date is the Psalter's bound.}
\midrule
Gradual & Ps 140:2 (Heb.\ 141) & David, by the title; an evening prayer beside
Israel's incense and evening sacrifice (v.~2) &
\chronodate{gradual}{\chronologyannotation{gradual}}\\
\cmidrule(lr){1-4}
\dossierprose{Titled \latin{Psalmus David}, ``A psalm of David'', again without
an occasion; the reign shown is the same reference point, not the psalm's date,
and the composition date is the Psalter's bound.}
\midrule
Gospel & Mt 22:1--14 & Writing: as Matthew left Palestine (Eusebius). Event:
Jerusalem, in the temple (21:23) &
\chronodate{gospel}{\chronologyannotation{gospel}}\\
\cmidrule(lr){1-4}
\dossierevent{Narrated event: the third parable Jesus speaks in the temple, on
the day after his entry into Jerusalem (21:1--18), when the chief priests and
ancients have questioned his authority (21:23), to the chief priests and
Pharisees (21:45); St Matthew gives it no date beyond that day.}
\cmidrule(lr){1-4}
\dossierprose{The event's year is derived: the \work{Catholic Encyclopedia}'s
``Jesus Christ'' (vol.~8, 1910), by A.~J. Maas, sets the parable in a journey
it dates by the founding of Rome, in a year it equates with A.D.~29. Traditional attribution: St Matthew the Apostle,
whom the Missal's heading names (\latin{secundum Matthaeum}). The
\work{Catholic Encyclopedia}'s ``Gospel of St.\ Matthew'' (vol.~10, 1911)
reports Eusebius's tradition that Matthew wrote his Gospel in Hebrew when he
left Palestine. Its five composition dates, shown after the event's year,
reckon from the Ascension, the Apostles' dispersal and St
Irenaeus, or give the Catholic critics' preference, and it settles on none, the
ancient writers being ``at variance''. The sixth, c.~A.D.~50, is Durand's in
``The New Testament'' (vol.~14, 1912) for the lost Aramaic original, of which
he adds, ``We know nothing definite of the date of its being rendered into
Greek.'' The comparison shown last is the modern critical horizon: the
\work{New American Bible, Revised Edition} places the Greek Gospel after
A.D.~70, a lower bound and not an exact date, which dates neither the Aramaic
original nor the parable's occasion.}
\midrule
Epistle & Eph 4:23--28 & Paul, ``a prisoner in the Lord'' (4:1), in Caesarea or
Rome; to ``the saints who are at Ephesus'' (1:1) &
\chronodate{epistle}{\chronologyannotation{epistle}}\\
\cmidrule(lr){1-4}
\dossierprose{Traditional attribution: St Paul, as the Missal's heading names
him, who names himself (1:1) and writes in chains (3:1; 6:20). The first range
shown is Ladeuze's in the \work{Catholic Encyclopedia}'s ``Epistle to the Ephesians''
(vol.~5, 1909), where of the letters of the captivity ``whether they were
produced in Caesarea or in Rome \dots\ is still a much mooted question''. The
year is from the table in Prat's ``St.\ Paul'' (vol.~11, 1911), which sets the
letter in the Roman captivity with Philemon, Colossians and Philippians; the
articles date that captivity differently, and neither answer is settled. The
comparison shown last is the modern critical horizon: the \work{New American
Bible, Revised Edition} dates around A.D.~80--100 the possibility that a later
disciple developed Paul's thought, which it offers beside a secretary working
at Paul's direction and the traditional captivity. That is a conditional
hypothesis, not a consensus, and it names no place for the later writing.}
\end{concisedossiertable}
\zlabel{triptych:concise:chronology:end}
